\documentclass[a4paper]{article}
\usepackage{graphicx}
\usepackage{amsmath,amssymb}
\usepackage{hyperref}
\usepackage{minted}
\usepackage{multirow}
\usepackage{float}
\usepackage{todonotes}
\usepackage{forest}

\title{}
\date{}

\begin{document}
    \maketitle
    Before attention is applied, the input values must first be transformed into numerical vectors. This step is called \textbf{encoding}. The purpose of encoding is simply to convert raw inputs into vectors that the model can process mathematically.

    For example, if the input is a time series
    
    \[
    x_1, x_2, x_3, x_4
    \]
    
    each element may originally be a scalar or a small feature vector. The encoder projects each element into a higher dimensional representation using a learned linear transformation:
    
    \[
    e_i = W_e x_i
    \]
    
    where \(W_e\) is a learned weight matrix. The resulting vector \(e_i\) is called the \textbf{embedding} of the input element.
    
    Thus the sequence becomes
    
    \[
    e_1, e_2, e_3, e_4
    \]
    
    These embeddings are then used to compute the Query, Key, and Value vectors used in attention.
    
    
    
    \section{}
    % \bibliographystyle{plain}
    % \bibliography{/home/donkarlo/Dropbox/repo/shared_research/bibliography/bibliography.bib}
\end{document}